% GENERATED FILE. Edit canonical lessons, then run npm run generate:books.
% canonical-source-hash: fnv1a64:b4771f7a0d0ebdb6
% canonical-lessons: KA-C530-virama, KA-C530-manaranjane, KA-C530-harate, KA-C530-asvadisu, KA-C530-maimare

\chapter{Free Time}
\label{ch:ka-free-time}
\hlchaptermodality{530}

\begin{chapteropening}
\textbf{By the end of this chapter:} \emph{I can say a pause, entertainment, a chat, to savour and to be absorbed.}

Complete the last lesson of chapter 530: I can say a pause, entertainment, a chat, to savour and to be absorbed.
\end{chapteropening}

\section[\texorpdfstring{\kn{ವಿರಾಮ} (virāma)}{ವಿರಾಮ (virāma)}]{\kn{ವಿರಾಮ} (virāma) — a pause; a break}

\label{lesson:KA-C530-virama}



\begin{quote}
Before the new one: say the Kannada for infatuation, then the Kannada for to live.
\end{quote}

\subsection*{You'll want to know: \kn{ವಿರಾಮ}}
\textbf{\kn{ವಿರಾಮ}} — \emph{virāma} — \textquotedblleft{}a pause; a break\textquotedblright{}.

\begin{quote}
\kn{ಎರಡು} \kn{ಗಂಟೆ} \kn{ಓದಿದ} \kn{ಮೇಲೆ} \kn{ನಾನು} \kn{ವಿರಾಮ} \kn{ತೆಗೆದುಕೊಂಡೆ}.
\end{quote}

\begin{quote}
\emph{eraḍu gaṇṭe ōdida mēle nānu virāma tegedukoṇḍe.}
\end{quote}

\begin{quote}
After studying for two hours I took a break.
\end{quote}

\kn{ವಿರಾಮದ} \kn{ವೇಳೆ}, leisure time. In grammar \kn{ಪೂರ್ಣವಿರಾಮ}, a full pause, is the full stop.

\begin{cousinweb}
\textbf{\kn{ವಿರಾಮ}} (\emph{virāma}) is borrowed from Sanskrit \emph{virāma}, \textquotedblleft{}cessation, termination\textquotedblright{} (Monier-Williams), from \emph{vi-ram}, \textquotedblleft{}to stop, to rest\textquotedblright{}.
\end{cousinweb}

\begin{grammarlens}[title={What you've built}]
One more word for describing what you lived and felt.
\end{grammarlens}

\subsection*{Your turn}
\begin{itemize}
\raggedright
  \item \emph{Say it:} \emph{virāma}
  \item \emph{Say it:} \emph{virāma}, once more
  \item \emph{Say it:} \emph{eraḍu gaṇṭe ōdida mēle nānu virāma tegedukoṇḍe.} — After studying for two hours I took a break.
  \item \emph{Recall:} say the Kannada for infatuation, then the Kannada for to live, then say \emph{virāma} again
\end{itemize}

\begin{tcolorbox}[breakable,colback=teal!4,colframe=teal!35!black,title={Before you move on}]
What does \kn{ವಿರಾಮ} mean? (\textquotedblleft{}A pause; a break\textquotedblright{}.) Say it once more.
\end{tcolorbox}
\section[\texorpdfstring{\kn{ಮನರಂಜನೆ} (manarañjane)}{ಮನರಂಜನೆ (manarañjane)}]{\kn{ಮನರಂಜನೆ} (manarañjane) — entertainment}

\label{lesson:KA-C530-manaranjane}



\begin{quote}
Before the new one: say the Kannada for to live, then the Kannada for a pause.
\end{quote}

\subsection*{You'll want to know: \kn{ಮನರಂಜನೆ}}
\textbf{\kn{ಮನರಂಜನೆ}} — \emph{manarañjane} — \textquotedblleft{}entertainment\textquotedblright{}.

\begin{quote}
\kn{ಟೀವಿ} \kn{ಅಜ್ಜಿಯ} \kn{ಮುಖ್ಯ} \kn{ಮನರಂಜನೆ}.
\end{quote}

\begin{quote}
\emph{ṭīvi ajjiya mukhya manarañjane.}
\end{quote}

\begin{quote}
The television is grandmother's main entertainment.
\end{quote}

\kn{ಮನ}, mind + \kn{ರಂಜನೆ}, delighting.

\begin{cousinweb}
\textbf{\kn{ಮನರಂಜನೆ}} (\emph{manarañjane}) is a learned compound built from Sanskrit \emph{manas}, \textquotedblleft{}mind\textquotedblright{}, and \emph{rañjana}, \textquotedblleft{}pleasing\textquotedblright{} (CDIAL 10587).
\end{cousinweb}

\begin{grammarlens}[title={What you've built}]
One more word for describing what you lived and felt.
\end{grammarlens}

\subsection*{Your turn}
\begin{itemize}
\raggedright
  \item \emph{Say it:} \emph{manarañjane}
  \item \emph{Say it:} \emph{manarañjane}, once more
  \item \emph{Say it:} \emph{ṭīvi ajjiya mukhya manarañjane.} — The television is grandmother's main entertainment.
  \item \emph{Recall:} say the Kannada for to live, then the Kannada for a pause, then say \emph{manarañjane} again
\end{itemize}

\begin{tcolorbox}[breakable,colback=teal!4,colframe=teal!35!black,title={Before you move on}]
What does \kn{ಮನರಂಜನೆ} mean? (\textquotedblleft{}Entertainment\textquotedblright{}.) Say it once more.
\end{tcolorbox}
\section[\texorpdfstring{\kn{ಹರಟೆ} (haraṭe)}{ಹರಟೆ (haraṭe)}]{\kn{ಹರಟೆ} (haraṭe) — a chat; idle talk}

\label{lesson:KA-C530-harate}



\begin{quote}
Before the new one: say the Kannada for a pause, then the Kannada for entertainment.
\end{quote}

\subsection*{You'll want to know: \kn{ಹರಟೆ}}
\textbf{\kn{ಹರಟೆ}} — \emph{haraṭe} — \textquotedblleft{}a chat; idle talk\textquotedblright{}.

\begin{quote}
\kn{ಸಂಜೆ} \kn{ಅಜ್ಜಂದಿರು} \kn{ಜಗಲಿಯ} \kn{ಮೇಲೆ} \kn{ಹರಟೆ} \kn{ಹೊಡೆಯುತ್ತಾರೆ}.
\end{quote}

\begin{quote}
\emph{sañje ajjandiru jagaliya mēle haraṭe hoḍeyuttāre.}
\end{quote}

\begin{quote}
In the evening the old men chat on the veranda.
\end{quote}

\kn{ಹರಟೆ} \kn{ಹೊಡೆ}, to chat away the time.

\begin{cousinweb}
\textbf{\kn{ಹರಟೆ}} (\emph{haraṭe}) is native Dravidian: DEDR 4031 lists Kannada \emph{haraṭu}, \textquotedblleft{}to talk idly\textquotedblright{}, and \emph{haraṭe}, \textquotedblleft{}idle talk\textquotedblright{}, beside Tamil \emph{paṟai}, \textquotedblleft{}to say\textquotedblright{}.
\end{cousinweb}

\begin{grammarlens}[title={What you've built}]
One more word for describing what you lived and felt.
\end{grammarlens}

\subsection*{Your turn}
\begin{itemize}
\raggedright
  \item \emph{Say it:} \emph{haraṭe}
  \item \emph{Say it:} \emph{haraṭe}, once more
  \item \emph{Say it:} \emph{sañje ajjandiru jagaliya mēle haraṭe hoḍeyuttāre.} — In the evening the old men chat on the veranda.
  \item \emph{Recall:} say the Kannada for a pause, then the Kannada for entertainment, then say \emph{haraṭe} again
\end{itemize}

\begin{tcolorbox}[breakable,colback=teal!4,colframe=teal!35!black,title={Before you move on}]
What does \kn{ಹರಟೆ} mean? (\textquotedblleft{}A chat; idle talk\textquotedblright{}.) Say it once more.
\end{tcolorbox}
\section[\texorpdfstring{\kn{ಆಸ್ವಾದಿಸು} (āsvādisu)}{ಆಸ್ವಾದಿಸು (āsvādisu)}]{\kn{ಆಸ್ವಾದಿಸು} (āsvādisu) — to savour; to relish}

\label{lesson:KA-C530-asvadisu}



\begin{quote}
Before the new one: say the Kannada for entertainment, then the Kannada for a chat.
\end{quote}

\subsection*{You'll want to know: \kn{ಆಸ್ವಾದಿಸು}}
\textbf{\kn{ಆಸ್ವಾದಿಸು}} — \emph{āsvādisu} — \textquotedblleft{}to savour; to relish\textquotedblright{}.

\begin{quote}
\kn{ಅಜ್ಜ} \kn{ಪ್ರತಿ} \kn{ಹಾಡನ್ನು} \kn{ಕಣ್ಣು} \kn{ಮುಚ್ಚಿ} \kn{ಆಸ್ವಾದಿಸುತ್ತಾರೆ}.
\end{quote}

\begin{quote}
\emph{ajja prati hāḍannu kaṇṇu mucci āsvādisuttāre.}
\end{quote}

\begin{quote}
Grandfather savours every song with his eyes closed.
\end{quote}

Past (I …): \kn{ಆಸ್ವಾದಿಸಿದೆ}, \emph{āsvādiside}.

Of food, music or a view.

\begin{cousinweb}
\textbf{\kn{ಆಸ್ವಾದಿಸು}} (\emph{āsvādisu}) is built on Sanskrit \emph{āsvāda}, \textquotedblleft{}eating with relish, tasting, enjoying\textquotedblright{} (Monier-Williams), with the verb ending -\emph{isu}.
\end{cousinweb}

\begin{grammarlens}[title={What you've built}]
One more word for describing what you lived and felt.
\end{grammarlens}

\subsection*{Your turn}
\begin{itemize}
\raggedright
  \item \emph{Say it:} \emph{āsvādisu}
  \item \emph{Say it:} \emph{āsvādisu}, once more
  \item \emph{Say it:} \emph{ajja prati hāḍannu kaṇṇu mucci āsvādisuttāre.} — Grandfather savours every song with his eyes closed.
  \item \emph{Recall:} say the Kannada for entertainment, then the Kannada for a chat, then say \emph{āsvādisu} again
\end{itemize}

\begin{tcolorbox}[breakable,colback=teal!4,colframe=teal!35!black,title={Before you move on}]
What does \kn{ಆಸ್ವಾದಿಸು} mean? (\textquotedblleft{}To savour; to relish\textquotedblright{}.) Say it once more.
\end{tcolorbox}
\section[\texorpdfstring{\kn{ಮೈಮರೆ} (maimare)}{ಮೈಮರೆ (maimare)}]{\kn{ಮೈಮರೆ} (maimare) — to be absorbed; to be lost in}

\label{lesson:KA-C530-maimare}



\begin{quote}
Before the new one: say the Kannada for a chat, then the Kannada for to savour.
\end{quote}

\subsection*{You'll want to know: \kn{ಮೈಮರೆ}}
\textbf{\kn{ಮೈಮರೆ}} — \emph{maimare} — \textquotedblleft{}to be absorbed; to be lost in\textquotedblright{}.

\begin{quote}
\kn{ಸಂಗೀತ} \kn{ಕೇಳುತ್ತಾ} \kn{ಅವಳು} \kn{ಮೈಮರೆತಳು}.
\end{quote}

\begin{quote}
\emph{saṅgīta kēḷuttā avaḷu maimaretaḷu.}
\end{quote}

\begin{quote}
Listening to the music, she lost herself in it.
\end{quote}

Past (I …): \kn{ಮೈಮರೆತೆ}, \emph{maimarete}.

\kn{ಮೈ}, the body + \kn{ಮರೆ}, to forget: to forget yourself.

\begin{cousinweb}
\textbf{\kn{ಮೈಮರೆ}} (\emph{maimare}) is a native Dravidian compound: \emph{mai}, \textquotedblleft{}the body\textquotedblright{} (DEDR 5073), and \emph{maṟe}, \textquotedblleft{}to forget\textquotedblright{} (DEDR 4760).
\end{cousinweb}

\begin{grammarlens}[title={What you've built}]
One more word for describing what you lived and felt.
\end{grammarlens}

\subsection*{Your turn}
\begin{itemize}
\raggedright
  \item \emph{Say it:} \emph{maimare}
  \item \emph{Say it:} \emph{maimare}, once more
  \item \emph{Say it:} \emph{saṅgīta kēḷuttā avaḷu maimaretaḷu.} — Listening to the music, she lost herself in it.
  \item \emph{Recall:} say the Kannada for a chat, then the Kannada for to savour, then say \emph{maimare} again
\end{itemize}

\begin{tcolorbox}[breakable,colback=teal!4,colframe=teal!35!black,title={Before you move on}]
What does \kn{ಮೈಮರೆ} mean? (\textquotedblleft{}To be absorbed; to be lost in\textquotedblright{}.) Say it once more.
\end{tcolorbox}
